\section{Periods, compact subgroups, and the transcendence question}

We use a Weierstrass equation over $\overline\Q$ and its invariant
differential. Its period lattice is $\Lambda=\Z\omega_1+\Z\omega_2$,
with $\Im(\omega_2/\omega_1)>0$. The quasi-periods are defined by
$\zeta(z+\omega_j)=\zeta(z)+\eta_j$. We extend
$\eta:\Lambda\to\C$ real-linearly to $\C$. The discriminant in
this section is the discriminant of this equation, equivalently of
the lattice in the specified differential. The quantity $N$ is
independent of changing this algebraic equation and differential.

\begin{proposition}[sigma expression and the third-kind period]\label{ari:C.7.1}
Let $P\in E(\overline\Q)$, with $P\ne O$ and $2P\ne O$, and
let $u$ be an elliptic logarithm of $P$. Then
\begin{equation}\label{ari:eq:cp-sigma}
 2\log N(P)=\frac16\log|\Delta|
       +\Re\{\Lambda_1+u(\zeta(u)-\eta(u))\},
 \qquad \Lambda_1=2\log\sigma(u)-u\zeta(u).
\end{equation}
The choices of logarithms change the expression inside the real
part by purely imaginary integral multiples of $2\pi$.
There is an extension $G_P$ of $E$ by $\mathbb G_m$, defined
over $\overline\Q$ and corresponding, under the principal
polarization, to $P$, for which $\Lambda_1$, modulo a logarithm
of an algebraic number, is the negative of a relative period
of $[\Z\to G_P]$. Changing an algebraic differential of the
third kind adds an algebraic multiple of $u$ and algebraic
endpoint terms.
\end{proposition}
\Uses{the hypotheses, notation and definitions stated here; \citet[Definitions 2.3--2.4 and
Section 6]{bertolin2022}; \cref{ari:eq:cp-periods}}

\begin{proof}
The sigma--theta product, or its divisor and automorphy law
together with its derivative at zero, gives
\begin{equation}\label{ari:eq:cp-klein}
 N(P)=|\Delta|^{1/12}|\sigma(u)|
                   \exp\!\left(-\frac12\Re(u\eta(u))\right).
\end{equation}
To check the constant, for $\Lambda=\Z+\tau\Z$ one uses
$\sigma(z)=e^{\eta_1z^2/2}\vartheta_1(z,\tau)/(2\pi\eta(\tau)^3)$,
$\Delta=(2\pi)^{12}\eta(\tau)^{24}$, and
$\eta(z)=\eta_1z-2\pi i\Im z/T$. Substitution gives exactly
the defining Gaussian in $N$. Homogeneity gives the general
lattice formula. Taking logarithms proves \eqref{ari:eq:cp-sigma}.

We make the period identification explicit. Put
\[
 f_u(z)=\frac{\sigma(z+u)}{\sigma(z)\sigma(u)}e^{-\zeta(u)z}.
\]
Its logarithmic differential is
\[
 \frac{df_u}{f_u}
 =\frac12\frac{\wp'(z)-\wp'(u)}{\wp(z)-\wp(u)}\,dz.
\]
This is the pullback of an algebraic differential with residues
$+1$ at $-P$ and $-1$ at $O$. The apparent $0/0$ at $z=u$
has a removable singularity when $P$ is not $2$-torsion.
The rational cocycle $f_u(z+z')/(f_u(z)f_u(z'))$ defines
the algebraic extension; in its rational trivialization the
exponential map has the form
\[
 \exp_{G_P}(z,v)=(\exp_E z,e^v f_u(z)).
\]
Its periods include
\begin{equation}\label{ari:eq:cp-periods}
 (0,2\pi i),\qquad
 (\omega_j,\zeta(u)\omega_j-\eta_j u),\quad j=1,2.
\end{equation}
These signs follow directly from the sigma automorphy law.
The identity
$\sigma(2u)=-\wp'(u)\sigma(u)^4$ implies
\[
 \log f_u(u)=\Lambda_1+\log(-\wp'(u)).
\]
For an algebraic lift $R=(P,r)$, $r\in\overline\Q^{\times}$,
the second coordinate of $\log_{G_P}R$ is
\begin{equation}\label{ari:eq:cp-v}
 v=\log r-\Lambda_1-\log(-\wp'(u)).
\end{equation}
It is a period of the relative homology defining the indicated
$1$-motive. This also avoids treating an integral with an endpoint
at a residue pole as an unregularized ordinary integral.
The description agrees with the generalized logarithm and
Serre function in \citet[Definitions 2.3--2.4 and
Section 6]{bertolin2022}; the signs in
\eqref{ari:eq:cp-periods} are for the exponential rather than the
quasi-periods of $f_u$ itself.
\end{proof}

\begin{proposition}[the compact-subgroup criterion and the real formula]\label{ari:C.7.2}
For $P$ as in Statement \ref{ari:C.7.1}, $N(P)$ is algebraic if and only
if the fibre of $G_P$ above $P$ contains an algebraic point of
the maximal compact subgroup $G_P(\C)^c$.
For the family in the stated conventions, if $t>1$ is real, choose
$\tau=iT$, $T>1$, and the zero $\mu=1/2-iy$, $0<y<T/2$.
With $x=2y/T$, one has
\[
 N(P)^4=3^{3/4}(t-1)^{-1/4}
       \left|\frac{\vartheta_1(x,i/T)}{\eta(i/T)}\right|.
\]
\end{proposition}
\Uses{the hypotheses, notation and definitions stated here; \cref{ari:C.7.1}; \cref{ari:eq:cp-periods}; \cref{ari:eq:cp-v}; \cref{ari:eq:cp-chi}; \citet[slides 5--8]{bertrand2017}}

\begin{proof}
On the analytic universal cover of $G_P$ define the real
homomorphism
\[
 \chi(z,v)=\Re\{v-\zeta(u)z+u\eta(z)\}.
\]
It vanishes on the periods \eqref{ari:eq:cp-periods}, hence descends
to $G_P(\C)$. Its restriction to the torus is $\log|\cdot|$.
The kernel is compact: in the universal cover it is a real
three-dimensional vector space modulo the rank-three period
lattice. Every compact subgroup has zero image in $\R$, so
this kernel is exactly $G_P(\C)^c$.
The identity
$\Re(u\eta(z))=\Re(z\eta(u))$ follows from the Legendre relation;
equivalently $\eta(z)=sz+\pi\bar z/\operatorname{area}(\Lambda)$
for a complex constant $s$. Substituting
\eqref{ari:eq:cp-v} at $z=u$ gives
\begin{equation}\label{ari:eq:cp-chi}
 \chi(R)=\log|r|-\log|\wp'(u)|
            -2\log N(P)+\frac16\log|\Delta|.
\end{equation}
If $r$ is algebraic and $\chi(R)=0$, this equation makes
$N(P)^2$ algebraic. Conversely, if $N(P)$ is algebraic,
choose the positive real algebraic number
$r=|\wp'(u)|\,|\Delta|^{-1/6}N(P)^2$.
Then \eqref{ari:eq:cp-chi} vanishes. This proves both directions
with the chosen rational trivialization of the fibre.
The compact interpretation is the one discussed by
\citet[slides 5--8]{bertrand2017}; here the calculation fixes
its normalization on the diagonal $P=P$.

For the real formula, the rectangular elliptic curve has its
two zeros of $\wp$ on $1/2+i\R$; along that edge $\wp$ runs
between the positive root $e_1$ and the negative root $e_2$.
Thus the indicated representative exists. Translate $2\mu$
by $-1$ and apply the modular substitution $\tau\mapsto-1/\tau$.
The resulting elliptic argument is $-2y/T$, which is real.
The Gaussian is therefore one. Evenness of $N$ and the
duplication identity prove the displayed formula. This removes
the Gaussian from this particular expression but makes no
algebraicity assertion about its theta value.
\end{proof}

\subsection{The conditional theorem and its exact conjectural input}

We use Grothendieck's period conjecture in the form stated as
Conjecture 1.1 of \citet{bertolin2022}: for a motive $M$ over
$\overline\Q$, the field generated by its periods over
$\overline\Q$ has transcendence degree $\dim G_{\mathrm{mot}}(M)$.
Only $1$-motives occur below. We also use their Construction 3.4
and Lemma 3.5: if $B$ is the smallest abelian subvariety
generated by the point and extension data, and $Z(1)$ is the
central torus constructed from the Weil-pairing bracket and
the remaining algebraic lift, then
\begin{equation}\label{ari:eq:cp-MTdim}
 \dim G_{\mathrm{mot}}(M)
 =\dim G_{\mathrm{mot}}(A)+2\dim B+\dim Z(1).
\end{equation}
The bracket is described explicitly in their Section 6.1.
Their Table 2 corrects the table in the paper they discuss; the
dependent nondeficient case has unipotent dimension three.
We now verify which case occurs here.

\begin{proposition}[the motivic groups needed on the diagonal]\label{ari:C.7.3}
Let $P$ be nontorsion and let $R$ be any algebraic lift above
$P$ in $G_P$. Put $M_R=[\Z\to G_P]$, $1\mapsto R$.
Its connected motivic group is, after a choice of Levi subgroup,
the semidirect product of $G_{\mathrm{mot}}(E)$ and the
three-dimensional Heisenberg group attached to $H_1(E)$ and
its principal polarization. Its dimension is $5$ for CM and
$7$ without CM.

For a non-real embedding put
$\mathcal M=M_R\oplus\overline{M_R}$.
If $E$ and $\bar E$ are not isogenous, its dimension is $13$.
If they are isogenous, choose an isogeny
$\phi:\bar E\to E$ of degree $d$ and write
$Q=\phi(\bar P)$. Let $r$ be the dimension over
$F=\End(E)\otimes\Q$ of the span of $P,Q$, modulo torsion.
Then
\[
\begin{array}{c|c|c|c}
\text{case}&\dim B&\dim Z&\dim G_{\mathrm{mot}}(\mathcal M)\\ \hline
r=2,\ E\text{ CM}&2&2&8\\
r=2,\ E\text{ not CM}&2&2&10\\
r=1,\ E\text{ CM}&1&1+e&5+e\\
r=1,\ E\text{ not CM}&1&1+e&7+e.
\end{array}
\]
Here $e$ is either zero or one: it is zero precisely when the
algebraic lift in the one-dimensional quotient of the two
central tori by the bracket torus is torsion. This criterion is
independent of the choices made modulo isogeny. In the real
case we use the single motive $M_R$; no extra quotient torus
is necessary.
\end{proposition}
\Uses{the hypotheses, notation and definitions stated here; \cref{ari:eq:cp-MTdim}}

\begin{proof}
All curves, polarizations, points, and lifts are algebraic over
$\overline\Q$, as required by Construction 3.4 and Lemma 3.5.
For $M_R$ the smallest abelian subvariety containing $(P,P)$ is
the diagonal $B=E$, since $P$ is nontorsion. If $\psi$ is the
polarization pairing, the bracket on this diagonal is
\[
 [(x,x),(y,y)]=\psi(x,y)-\psi(y,x)=2\psi(x,y).
\]
It is nonzero and spans $\Q(1)$. Hence $Z'=Z=\mathbb G_m$;
there is no deficient or Ribet exception, including for CM.
The pure group is $\mathrm{GL}_2$ without CM and
$\operatorname{Res}_{F/\Q}\mathbb G_m$ for the imaginary
quadratic CM field. Their dimensions are four and two.
This proves the single-motive assertion, including its Lie
bracket and Levi action: the action on the centre is the
polarization multiplier.

For the direct sum only the two diagonal central coordinates
occur. If the two elliptic curves are not isogenous, they
cannot be CM conjugates. The pure group is
$\{(g_1,g_2)\in\mathrm{GL}_2^2:\det g_1=\det g_2\}$,
of dimension seven. The point data generate $E\times\bar E$,
and the two brackets are independent, giving $7+4+2=13$.

If the curves are isogenous, transport the second polarized
factor through $\phi$. Its polarization pairing becomes
$\psi/d$. The smallest abelian subvariety has dimension $r$.
For $r=2$, independent variation in its two factors makes the
two diagonal brackets span both central coordinates; thus
$Z'=Z=\mathbb G_m^2$. The unipotent group is the product
of two Heisenberg groups, with the same pure group acting
on both standard factors.
For $r=1$, after multiplying the points by a common integer
we may write $Q=aP$, $a\in F^\times$. The bracket image is
the one-dimensional line with proportional central coordinates
\begin{equation}\label{ari:eq:cp-positive-centre}
 (1,c),\qquad c=\frac{a a^\dagger}{d}>0,
\end{equation}
where $\dagger$ is the Rosati involution. On an elliptic curve
$aa^\dagger$ is a positive rational number: a square in the
non-CM case and the positive field norm in the CM case.
Construction 3.4(3) adds a further central direction exactly
when the image of the algebraic lift in the quotient torus
is nontorsion. This is the integer $e$ above.
The unipotent Lie algebra is the Heisenberg algebra, with one
additional central Tate line when $e=1$. Formula
\eqref{ari:eq:cp-MTdim} proves every entry of the table. This
description computes the group even when the lift contributes
a Kummer factor; specifying only the span of $P,Q$ would
not determine that last factor.
\end{proof}

\begin{proposition}[conditionally: transcendence of the local norm]\label{ari:C.7.4}
Assume Grothendieck's period conjecture for $1$-motives over
$\overline\Q$. For every algebraic parameter $t\ne0,1$ at
which $P_t$ is nontorsion, $N(P_t)$ is transcendental.
More generally the conclusion holds for every nontorsion
algebraic point on an elliptic curve over $\overline\Q$.
The weaker named Relative Semi-abelian Conjecture 4.7 of
\citet{bertolin2022}, restricted to the groups and algebraic
lifts used in this proof, already suffices.
\end{proposition}
\Uses{the hypotheses, notation and definitions stated here; \citet{bertolin2022}; \cref{ari:C.7.2}; \cref{ari:eq:cp-v}; \citet[Remark 4.13]{bertolin2022}; \cref{ari:eq:cp-positive-centre}; \cref{ari:eq:cp-relation}; \cref{ari:eq:cp-real-linear}}

\begin{proof}
Suppose $N(P)$ is algebraic. Statement \ref{ari:C.7.2} supplies an
algebraic lift $R=(P,r)$ in $G_P(\C)^c$, and its conjugate
supplies a lift $\bar R$ in $G_{\bar P}(\C)^c$.
Let $v$ denote the logarithmic coordinate in
\eqref{ari:eq:cp-v}. Put
\[
 K_0=\overline\Q\bigl(\text{entries of the full de Rham--Betti
                    period matrix of }G_P\times G_{\bar P}\bigr).
\]
Here ``full'' includes the integrals of differentials of the second
kind, hence the quasi-periods $\eta_j$, as well as the periods
$\omega_j$, the third-kind periods $\kappa_j$, and $2\pi i$.
It is the matrix $\Omega_G$ in their Definition 2.4, not merely
the lattice in the Lie algebra of the semi-abelian variety.
This field contains $u,\zeta(u)$ and their complex conjugates.
Indeed it contains the pure periods, $2\pi i$, and
$\kappa_j=\zeta(u)\omega_j-\eta_j u$; the determinant of
the two equations for $u,\zeta(u)$ is, up to sign,
$\omega_1\eta_2-\omega_2\eta_1= -2\pi i$.
This is also the observation of \citet[Remark 4.13]{bertolin2022}.
The real-linear expression for $\eta(u)$ is
\begin{equation}\label{ari:eq:cp-real-linear}
 \eta(u)=
 \frac{(u\bar\omega_2-\bar u\omega_2)\eta_1
       +(\bar u\omega_1-u\bar\omega_1)\eta_2}
      {\omega_1\bar\omega_2-\bar\omega_1\omega_2}.
\end{equation}
The denominator is nonzero. Thus $u\eta(u)$ and its conjugate
belong to $K_0$. The compact equation $\chi(R)=0$ becomes
the algebraic relation over $K_0$
\begin{equation}\label{ari:eq:cp-relation}
 v+\bar v
 =u\zeta(u)+\bar u\overline{\zeta(u)}
                -u\eta(u)-\bar u\overline{\eta(u)}.
\end{equation}
Complex conjugates here are coordinates of the conjugate
$1$-motive. They are not being treated as rational functions
of un-conjugated periods without justification.

Consider the period torsor of $\mathcal M$ over the period
torsor of $G=G_P\times G_{\bar P}$. The kernel of the
corresponding motivic groups is its weight $-2$ vector group
$\operatorname{Lie}Z$, of dimension $\dim Z$. There is no
weight $-1$ kernel: the point coordinates are exactly the
extension coordinates under the polarization, so $B_Q=0$.
More explicitly, projection of the smallest abelian variety $B$
to the variety generated by the extension coordinates is an
isogeny, since $B$ is the graph of the same point data in the
two factors. It therefore induces an isomorphism on the
weight $-1$ graded Lie algebras. The motive $[0\to G]$ has
only weights $-1$ and $-2$, and its unipotent Lie algebra has
no weight $-2$ graded piece. Thus the entire kernel is exactly
$W_{-2}\operatorname{Lie}G_{\rm mot}(\mathcal M)$; no additional
abelian direction or unaccounted logarithm remains in the fibre.
After evaluating the base periods, the fibre is an affine
torsor over $K_0$ under that vector group. Translation in
its two central coordinates changes $(v,\bar v)$ by
$2\pi i(c_1,c_2)$, with the signs fixed by using a positive
Tate loop for each motive. The linear functional
$c_1+c_2$ is nonzero on this central group. It is nonzero
on the whole plane when $\dim B=2$, and when $\dim B=1$
it is already nonzero on the bracket line
\eqref{ari:eq:cp-positive-centre}, because $1+c>0$.

Consequently \eqref{ari:eq:cp-relation} cuts a proper affine
hyperplane in the period-torsor fibre. It forces the relative
transcendence degree of the actual periods over $K_0$ to be
at most $\dim Z-1$. The pure-base period torsor has dimension
$\dim G_{\mathrm{mot}}(G)$, so the total transcendence degree
is at most $\dim G_{\mathrm{mot}}(\mathcal M)-1$.
This contradicts the asserted period conjecture. For real
data the single motive suffices: conjugation acts on its
Betti realization over $\Q$, and the same equation has a
nonzero coefficient on its one central direction. One can
also double the point and lift into the identity component,
choose $u,\omega_1$ real, and read the relation directly as
\[
 \omega_1 v-\omega_1u\zeta(u)+\eta_1u^2=0.
\]
Doubling preserves the equivalence of algebraicity of $N$
by the division formula and its algebraic nonzero factors.

For the weaker conjectural input, Conjecture 4.7 asserts
relative transcendence degree at least $2\dim B_Q+\dim Z$.
All $R_i$ used here are algebraic, $B_Q=0$, and the
generalized logarithms add exactly the central coordinates
to $K_0$. Thus it contradicts the same hyperplane relation
directly. Proposition 4.12 of that paper derives Conjecture
4.7 from the generalized period conjecture; in our case the
field of definition is algebraic, so the ordinary period
conjecture is enough. The weaker Conjecture 4.9, formulated
using arbitrary algebraic translates, gives a zero lower
bound at an algebraic point and does not by itself yield
this conclusion. Linear relations among $1$-periods are
also insufficient: after clearing the denominator in
\eqref{ari:eq:cp-real-linear}, the prohibited relation is
polynomial, not a linear relation with algebraic coefficients.
\end{proof}


\paragraph{Cited input.} \label{ari:C.7.5}The source statement and its expanded derivation are retained in the proof compendium, Paper~I, label \nolinkurl{ari:C.7.5}.\par




\subsection{CM theta normalization and further unconditional deductions}

The holomorphic theta normalization proposed for the CM interpolation approach differs from the
norm occurring in the stated problem. We make that distinction explicit,
including the finite-character ambiguity in the multiplication law.

\begin{proposition}[the exact CM multiplication law]\label{ari:C.7.6}
Let $E/\overline{\Q}$ have complex multiplication by an order
$\mathcal O$, and use an algebraic invariant differential. Write
$A=\operatorname{area}(\Lambda)$ and
\[
 \eta(z)=s_2z+\frac{\pi}{A}\bar z,\qquad
 \theta(z)=\sigma(z)e^{-s_2z^2/2}.
\]
Then $s_2\in\overline{\Q}$. For every nonzero
$\alpha\in\mathcal O$, with $d=\alpha\bar\alpha$, there is
$F_\alpha\in\overline{\Q}(E)$ such that
\begin{equation}\label{ari:eq:ct-CMmult}
 \theta(\alpha z)^2=F_\alpha(\exp_E z)\theta(z)^{2d}.
\end{equation}
Its divisor and leading Laurent coefficient are
\[
 \operatorname{div}(F_\alpha)=2[\alpha]^*(O)-2d(O),\qquad
 F_\alpha(\exp_Ez)=\alpha^2z^{2-2d}(1+O(z^2)).
\]
Thus the unsquared quotient is an algebraic function on $E$, but
it need not be a rational function on $E$. For a nontorsion
$P=\exp_Eu\in E(\overline{\Q})$, algebraicity of $\theta(u)$
would imply algebraicity of $\theta(\alpha u)$ for all
$\alpha\in\mathcal O$. The additive group $\mathcal O u$ has
rank two, not more than four.
\end{proposition}
\Uses{the hypotheses, notation and definitions stated here; the explicit constructions and calculations in the proof}

\begin{proof}
The Legendre relation gives the displayed real-linear decomposition
of $\eta$. The sigma automorphy formula gives
\[
 \theta(z+\omega)=\epsilon(\omega)
  \exp\left(\frac{\pi\bar\omega}{A}
                         (z+\omega/2)\right)\theta(z),
 \qquad \epsilon(\omega)\in\{1,-1\}.
\]
For $\alpha\Lambda\subset\Lambda$, the exponential multipliers
of $\theta(\alpha z)/\theta(z)^d$ cancel because
$\alpha\bar\alpha=d$. Its remaining multiplier is a sign.
Squaring therefore gives a meromorphic function on $E$, with the
stated divisor. This divisor is defined over $\overline{\Q}$.
The existence of a rational function with this divisor over $\C$
implies its existence over $\overline{\Q}$: the relevant
one-dimensional Riemann--Roch space is obtained by extension of
scalars. Normalize its leading Laurent coefficient, in the formal
parameter of the algebraic differential, to $\alpha^2$. This
identifies it with the squared analytic quotient. The action of
$\alpha$ on this differential is algebraic, so $\alpha^2$ is an
allowed normalization.

Choose a nonreal $\alpha\in\mathcal O$. Since the normalized
sigma series is $\sigma(z)=z+O(z^5)$ with algebraic Laurent
coefficients, the next coefficient of the squared quotient is
\[
 \alpha^2z^{2-2d}
       \{1-s_2(\alpha^2-d)z^2+O(z^4)\}.
\]
The coefficients of the rational function in this algebraic formal
parameter are algebraic. As $\alpha^2-d\ne0$, this proves
$s_2\in\overline{\Q}$ without an unnormalized appeal to a CM
weight-two form. It also proves that the functions in
\eqref{ari:eq:ct-CMmult} are algebraic over the chosen field.
At nontorsion $u$, none of the arguments $\alpha u$ is a lattice
point; evaluating the identity proves the algebraicity implication.
The injective map $\mathcal O\to\C$, $\alpha\mapsto\alpha u$,
proves the rank assertion.

The qualification about rationality is necessary. For a square
lattice, $s_2=0$ and $\alpha=1+i$ give
$F_{1+i}=2i\wp$. The unsquared quotient has a simple zero at
the nonzero point in $\ker(1+i)$ and a simple pole at $O$,
which is not a principal divisor on $E$; its square is the
rational function just displayed.
\end{proof}


\paragraph{Cited input.} \label{ari:C.7.7}The source statement and its expanded derivation are retained in the proof compendium, Paper~I, label \nolinkurl{ari:C.7.7}.\par

\Uses{the hypotheses, notation and definitions stated here; \cref{ari:C.7.6}; \citet[Theorem~1 and Corollary~1]{waldschmidt1979sigma}; \citet[Theorem~6.1]{huber2022}; \cref{app:rederive:ari-C.7.7}; \citet[final slide, comment (c)]{bertrand2017}; \citet[p.~15-03]{waldschmidt1988theta}; \cref{ari:C.7.1}; \cref{ari:eq:ct-Gaussian}; \citet[Corollary~2]{reyssat1982weierstrass}; \citet[p.~15-05]{waldschmidt1988theta}; \citet[Theorem 3.1.2]{bost2009}; \cref{ari:C.7.2}; \citet[Theorem 1.4 and Corollary 1.6]{nguyen2023}; \citet[Theorems 1.3 and 2.3]{xie2024algebraicity}; \cref{ari:C.7.9}}


The stated normalized values follow from \citet[Theorem~1 and Corollary~1]{waldschmidt1979sigma} and the analytic subgroup theorem in \citet[Theorem~6.1]{huber2022}. A complete re-derivation is given in Appendix~\ref{app:rederive:ari-C.7.7}.


\paragraph{A theta criterion requiring its original normalization.}
An earlier draft applied the pointwise criterion printed in
\citet[final slide, comment (c)]{bertrand2017} to
$g(u,v)-s_2uv$, where
$g(u,v)=\log(\sigma(u+v)/(\sigma(u)\sigma(v)))$.
It inferred transcendence of $\theta(u)$ by taking $v=u$.
We withdraw that inference pending verification of the original
lift conventions in \citet[Corollary 3]{bertrand1984}.
Indeed replacing $v$ by $v+\omega$, $\omega\in\Lambda$,
changes $g(u,v)-s_2uv$, modulo $2\pi i\Z$, by
\[
 (\eta(\omega)-s_2\omega)u=\frac{\pi\bar\omega u}{A},
\]
whereas the condition $N\exp_Ev=f(\exp_Eu)$ with $f$
Rosati-antisymmetric depends only on the elliptic points.
Thus the two-sided statement cannot be used with unrestricted
lifts without checking its normalization. The author's slide
and the original paper's first two pages were accessible;
the full original Corollary 3 was not. We use neither direction
of that criterion here. The withdrawn text and the lift audit
are retained in the records.

A subsequent source check resolves the normalization of the
one-variable holomorphic assertion itself:
\citet[p.~15-03]{waldschmidt1988theta} explicitly states that
$\sigma(u)e^{-s_2u^2/2}$ is transcendental when the invariants
are algebraic, $E$ has CM, and $\exp_Eu$ is an algebraic
nontorsion point.  That source defines $s_2$ by the regularized
lattice sum and attributes the result to Bertrand's 1982
paper \emph{Valeurs de fonctions th\^eta et hauteurs p-adiques}.
We have checked this complete normalized statement in
Waldschmidt's published exposition; the full original 1982
proof has not been retrieved.  This source clarification
does not validate the unrestricted-lift inference withdrawn
above, and the theorem concerns the holomorphic value rather
than the metric norm.  \ref{ari:C.7.7} retains the two assertions whose
proofs have been fully checked here.

\paragraph{The archimedean term still requiring control.}
The exact comparison is
\begin{equation}\label{ari:eq:ct-Gaussian}
 N(P)=|\Delta|^{1/12}|\theta(u)|
                         \exp\left(-\frac{\pi|u|^2}{2A}\right).
\end{equation}
It follows from \ref{ari:C.7.1} and the real-linear decomposition of
$\eta$. The normalization proved transcendental in \ref{ari:C.7.7} instead
gives the exact identity
\[
 N(P)=|\Delta|^{1/12}
 \left|\sigma(u)e^{-u\zeta(u)/2}\right|
 \exp\left(\frac12\Re\{u(\zeta(u)-\eta(u))\}\right).
\]
Neither formula excludes algebraicity of $N(P)$: transcendence
of one complex factor does not control its absolute value or
its product with the remaining exponential.
The generic finite-order interpolation approach
supplies only the rank-two set $\mathcal O u$; it does not
meet the proposed rank-greater-than-four criterion. The
derivative approach would need the algebraicity of
$\theta'(u)/\theta(u)=\zeta(u)-s_2u$, which \ref{ari:C.7.7} excludes.
Replacing the norm by the holomorphic theta value would
therefore change the problem precisely by the factor in
\eqref{ari:eq:ct-Gaussian}.

The related independence theorem of
\citet[Corollary~2]{reyssat1982weierstrass} says that two of
$e^{i\pi u/\omega_1}$, $s(u)$ and $s'(u)$ are algebraically
independent, where $s(z)=\sigma(z)e^{-\eta_1z^2/(2\omega_1)}$.
For real $u$, hypothetical algebraicity of $N(P)$ makes
$s(u)$ algebraic, so the conclusion instead forces independence
of $e^{i\pi u/\omega_1}$ and
$\zeta(u)-\eta_1u/\omega_1$; it yields no contradiction.
The periodic-functional-equation theorem reported in
\citet[p.~15-05]{waldschmidt1988theta} assumes algebraic
$q=e^{i\pi\tau}$, whereas this $q$ is transcendental for
imaginary quadratic $\tau$ by Gelfond--Schneider.  These
specific checked applications therefore do not decide the
metric norm; their source scopes and the larger CM
independence tuples are recorded separately.

We also tested an additive-extension route.
\citet[Theorem 3.1.2]{bost2009} proves that an algebraic point
of the universal vector extension of an abelian variety which
lies in its maximal compact subgroup at one complex embedding
is torsion. The theorem's hypotheses require an algebraic
point of the vector extension, equivalently an algebraic line
bundle with an algebraic connection having unitary monodromy.
\ref{ari:C.7.2} supplies only an algebraic point in one fibre of a
multiplicative extension. It does not supply that algebraic
connection. Differentiating a metric value at one algebraic
point cannot establish algebraicity of the connection there
or globally. This is the exact missing implication in that
attempt.

Another route uses the Latt\`es map $f_E$ defined by
$x(2Q)=f_E(x(Q))$. Its canonical height is
$\widehat h_{f_E}(x(P))=2\widehat h_O(P)$. For rational
parameters C.6 therefore makes algebraicity of $N(P)$
equivalent to algebraicity of
$\exp(\widehat h_{f_E}(x(P)))$.
\citet[Theorem 1.4 and Corollary 1.6]{nguyen2023} settle the
corresponding multiplicative-height problem for nonintegrable
polynomials using B\"ottcher coordinates in the basin of
infinity. They do not apply to $f_E$. Indeed an elliptic
duplication Latt\`es map has no exceptional point: the inverse
images of any point contain the images of all successive
division points of an elliptic point, an infinite set. A
polynomial has the exceptional point $\infty$, so no conjugate
of $f_E$ is polynomial. The Julia set is the whole sphere
and the required basin of infinity is absent. This checks
the failed hypothesis, rather than treating a theorem for
polynomials as a theorem for all rational maps.

We also checked the arithmetic algebraicity criterion of
\citet[Theorems 1.3 and 2.3]{xie2024algebraicity}.
Besides sufficiently rapid height-normalized approximation by
points in a fixed number field, it requires the analytic germ
to be formally defined over that field. This hypothesis fails
for the compact leaf in \ref{ari:C.7.9}: its tangent has the proved
transcendental slope
$c=\zeta(u)-\eta_1u/\omega_1$.
A smooth formal germ defined over $\overline{\Q}$ has an
algebraic tangent; algebraic translations, coordinate changes
and isogenies preserve this obstruction. Assuming $N(P)$
algebraic supplies an algebraic point on the leaf and its
integer multiples, but not an algebraic formal germ. Attempting
to force convergence by successive division also loses the
fixed-number-field hypothesis: if $[2]^nQ_n=Q$ with all
$Q_n\in E(K)$ for one number field, then
$\widehat h_O(Q_n)=4^{-n}\widehat h_O(Q)>0$.
The $Q_n$ are distinct with bounded ample height, contrary to
Northcott finiteness. The full hypothesis audit and the opened
preprint are retained in the records.

\paragraph{A current source whose theorem text remains unchecked.}
The publisher's abstract for \citet{pham2025sigma} announces
complex and $p$-adic linear-independence results for semiabelian
varieties and applications to Weierstrass sigma functions.
We opened its metadata, abstract and reference list, but both
the full-HTML link and the publisher's citation-PDF link returned
the abstract/access page. No accessible author manuscript was
located. Consequently we could not check any theorem statement
or its hypotheses. This is a potentially relevant unverified
source: we make no claim either that it settles the Gaussian
norm question or that its results are inapplicable to it.



\subsection{The precise nonarchimedean comparison}

\paragraph{Cited input.} \label{ari:C.7.8}The source statement and its expanded derivation are retained in the proof compendium, Paper~I, label \nolinkurl{ari:C.7.8}.\par

\Uses{the hypotheses, notation and definitions stated here; \citet[Appendix A]{agboola2006}; \citet{agboola2006}; \cref{app:rederive:ari-C.7.8}; \cref{ari:C.7.6}; \citet[Theorem 2.1]{bertrand1995}; \cref{ari:C.7.3}; \cref{ari:C.7.4}; \citet[Theorem 3.1]{bertrand1995}; \citet{bost2009}; \cref{ari:C.7.2}}


The nonvanishing theorem is reproduced in Rubin's Appendix~A, Theorem~A.1 and Lemma~A.2 of \citet{agboola2006}; the exact hypotheses and application are retained. A complete re-derivation is given in Appendix~\ref{app:rederive:ari-C.7.8}.


The comparison explains the missing step in an attempted
archimedean adaptation. The ordinary $p$-adic sigma construction
uses a holomorphic quadratic correction. In the CM normalization
the corresponding complex holomorphic function is $\theta$
of \ref{ari:C.7.6}, whereas the real metric also contains
$\exp(-\pi|u|^2/(2A))$. Theorem A.1 supplies no identity
between the ordinary $p$-adic height and that real factor.
Transferring the conclusion would require precisely such a
new argument; replacing the real factor by the holomorphic
CM correction does not preserve the norm being asked about.

There is also a definite limitation on the maximal-compact
approach. \citet[Theorem 2.1]{bertrand1995} constructs
nontrivial multiplicative extensions of polarized abelian
varieties with nontorsion algebraic points in their maximal
real compact subgroup. Its hypotheses include an endomorphism
not fixed by the Rosati involution and a nontorsion rational
point over a real number field. The construction uses a base
point related to the extension class by a Rosati-antisymmetric
endomorphism. It cannot be replaced by a theorem saying that
every algebraic compact point of every multiplicative
extension is torsion. For our diagonal class, an equality
$NP=f(P)$ with $N>0$ and $f+f^\dagger=0$ would force
$P$ torsion: $N-f$ has no zero eigenvalue, since
$(N-f)^\dagger(N-f)=N^2-f^2$ is positive. In the elliptic
case this is simply a nonzero isogeny. The diagonal positivity
used in \ref{ari:C.7.3}--\ref{ari:C.7.4} is therefore essential, but these
counterexamples neither prove nor refute the desired
unconditional diagonal statement. The additive result in
\citet[Theorem 3.1]{bertrand1995}, and its higher-dimensional
extension in \citet{bost2009}, requires a point in a vector
extension, a hypothesis not furnished by \ref{ari:C.7.2}.


\subsection{The fixed semiabelian extension and the diagonal subgroup attempt}

\begin{proposition}[the endomorphism and slope obstructions]\label{ari:C.7.9}
Let $E/\overline\Q$ be an elliptic curve, let
$q\in E^\vee(\overline\Q)$ be nontorsion, and let
\[
 1\longrightarrow\mathbb G_m\longrightarrow G_q
       \longrightarrow E\longrightarrow0
\]
be the extension of class $q$. Then
$\operatorname{End}_{\overline\Q}(G_q)=\Z$, including when
$E$ has complex multiplication. In particular a noninteger
CM endomorphism of $E$ does not lift to an endomorphism of this
fixed extension.

For
\[
 E:y^2+y=x^3-38x+90,\qquad P=(0,9),\qquad Q=2P=(4,-2),
\]
one has
\[
 N(Q)=19^{1/4}N(P)^4.
\]
Use the differential $dx/(2y+1)$, choose the positive real
period $\omega_1$, and choose the real logarithm
$u\in(0,\omega_1/2)$ of $Q$. Put
\begin{equation}\label{ari:eq:cd-slope}
 c=\zeta(u)-\frac{\eta_1}{\omega_1}u,
 \qquad \eta_1=\eta(\omega_1).
\end{equation}
Then $c$ is a real transcendental number. If $N(P)$ were
algebraic, the extension $G_Q$ of \ref{ari:C.7.1} would have a
nontorsion algebraic point whose logarithm is $(u,cu)$.
The complex line
$\ell_c=\{(z,cz):z\in\C\}$ is not defined over
$\overline\Q$ in the algebraic logarithm coordinates of
\ref{ari:C.7.1}. Thus the analytic subgroup theorem cannot be applied
to that line using the hypothesized compact point.
\end{proposition}
\Uses{the hypotheses, notation and definitions stated here; \cref{ari:C.7.1}; \cref{ari:eq:cd-compatibility}; \citep[Theorem 6.1]{huber2022}; \cref{ari:C.6.R1}; \cref{ari:C.7.2}}

\begin{proof}
The torus is the maximal connected affine subgroup of $G_q$.
Consequently every endomorphism of $G_q$ induces multiplication
by some $n\in\Z$ on $\mathbb G_m$ and an endomorphism
$\alpha$ on $E$. Functoriality of extension classes gives
\begin{equation}\label{ari:eq:cd-compatibility}
 \alpha^*q=nq\quad\text{in }E^\vee.
\end{equation}
Here $\alpha^*$ is the dual endomorphism, not $\alpha$
acting on an arbitrarily identified copy of $E$. To make the
convention explicit, choose a principal polarization
$\lambda:E\to E^\vee$, put $Q_0=\lambda^{-1}q$, and
write $\alpha^\dagger=\lambda^{-1}\alpha^*\lambda$.
Equation \eqref{ari:eq:cd-compatibility} says
$(\alpha^\dagger-[n])Q_0=O$. A nonzero endomorphism of an
elliptic curve is an isogeny and has finite kernel. Since
$Q_0$ is nontorsion, $\alpha^\dagger=[n]$, hence
$\alpha=[n]$ as well.

The difference between the given endomorphism and $[n]$ has
image in the torus and is trivial on the torus. It therefore
factors through an algebraic-group homomorphism
$E\to\mathbb G_m$. This homomorphism is zero, since a
morphism from the proper connected curve $E$ to the affine
group $\mathbb G_m$ is constant and a group homomorphism
preserves the identity. Thus there is no additional ambiguity
from $\operatorname{Hom}(E,\mathbb G_m)$, and the
endomorphism is exactly $[n]$. Conversely all integer
multiplications exist. This proves the first assertion.

For the concrete curve, the tangent slope at $P$ is $-2$,
giving $2P=(4,-2)$. Its discriminant is $-19^3$ and
$\psi_2(P)=2\cdot9+1=19$. The division identity in C.4
therefore gives the asserted formula for $N(Q)$, so replacing
$P$ by $Q$ preserves algebraicity of the norm. This doubling
also places the base point in the identity component of the
real elliptic group without any component assumption. In the
chosen differential the Weierstrass coordinates are
$(\wp,\wp')=(x,2y+1)$, with $g_2=152$, $g_3=-361$.
The negative discriminant means that the real cubic has one
real root. Along $(0,\omega_1/2)$, $\wp$ decreases from
$+\infty$ to that root and $\wp'<0$. Since
$\wp'(u)=2(-2)+1=-3$, the indicated real logarithm of $Q$
exists.

The period vectors of $G_Q$ in \ref{ari:C.7.1} include
\[
 (\omega_1,\zeta(u)\omega_1-\eta_1u)
       =\omega_1(1,c).
\]
These are algebraic logarithm coordinates: $dz$ and the
invariant differential $d\log r-df_u/f_u$ are algebraic,
because $df_u/f_u$ is the algebraic third-kind differential
displayed in \ref{ari:C.7.1}. Real invariants and real $u,\omega_1$
make $c$ real.

Suppose $c$ were algebraic. Then $\ell_c$ would be an
algebraic Lie subspace. Dividing the nonzero period
$\omega_1(1,c)$ by a sufficiently large positive integer
gives a vector in $\ell_c$ whose exponential is a nonzero
torsion point of $G_Q$, hence an algebraic point. The
analytic subgroup theorem
\citep[Theorem 6.1]{huber2022} supplies a positive-dimensional
connected algebraic subgroup with Lie algebra contained in
$\ell_c$. Its dimension is one, and its projection to $E$
is nontrivial. Such a subgroup must be an elliptic curve
$E'$ projecting isogenously to $E$: a connected
one-dimensional affine algebraic group cannot map
nontrivially to an elliptic curve. Its inclusion splits the
pullback of $G_Q$ to $E'$. The dual of that isogeny would
therefore kill the extension class $\lambda(Q)$, forcing
$Q$ torsion. This contradicts \ref{ari:C.6.R1}. We conclude that $c$
is transcendental. In particular the line is not defined over
$\overline\Q$. This argument uses the period vector alone;
it does not assume algebraicity of $N(P)$.

Finally assume, only for the attempted contradiction, that
$N(P)$ is algebraic. The positive real number
\[
 r=3\,19^{-1/2}N(Q)^2
\]
is algebraic, and \ref{ari:C.7.2} places the point $R=(Q,r)$ of
$G_Q$ in its maximal compact subgroup. The chosen real
logarithm has
\[
 f_u(u)=-\wp'(u)\sigma(u)^2e^{-u\zeta(u)}
        =3\sigma(u)^2e^{-u\zeta(u)}>0.
\]
Consequently $v=\log(r/f_u(u))$ is real and
$(u,v)$ is a logarithm of $R$, with no imaginary logarithm
ambiguity. Substituting this real logarithm into the exact
compactness equation $\chi=0$ in \ref{ari:C.7.2} gives
\[
 v=\zeta(u)u-u\eta(u)
   =u\left(\zeta(u)-\frac{\eta_1}{\omega_1}u\right)=cu.
\]
Its projection $Q$ is nontorsion, so $R$ is nontorsion.
The input missing from this proposed use of the analytic
subgroup theorem is a proper Lie subspace defined over
$\overline\Q$; the actual compact slope line is proved
not to have that property. The CM decomposition makes the
same obstruction visible as
\[
 c=\zeta(u)-s_2u-\frac{\pi}{A}u.
\]
The last term is the derivative of the real Gaussian
correction. It cannot be discarded while retaining the
compactness equation.
\end{proof}

\paragraph{What functional transcendence gives in this construction.}
The functional theorem can be applied fully, but its conclusion
does not contradict this hypothetical algebraic compact point.
Indeed consider
\[
 f:\C\longrightarrow\operatorname{Lie}(G_Q),\qquad
 f(z)=(z,cz).
\]
No proper connected algebraic subgroup of $G_Q$ projects
onto $E$: as in the proof above, such a subgroup would split
the nontorsion extension after an isogeny. Every proper
connected subgroup is therefore contained in the torus,
whereas $\exp_{G_Q}f$ has nonconstant elliptic projection.
Thus $f$ is $G_Q$-nondegenerate in Definition 1.2 of
\citet{noguchi2024}. Their Theorem 1.3 gives
\[
 \dim_{\C}\overline{\{(\exp_{G_Q}f(z),f(z)):z\in\C\}}^{\rm Zar}
       \ge \dim G_Q+1=3.
\]
The graph is contained in $G_Q\times\ell_c$, of dimension
three, so equality holds. This is a statement over $\C$
about a family of meromorphic functions. It remains valid
whether or not the single value $\exp_{G_Q}f(u)$ is
algebraic. Under the norm-algebraicity hypothesis the integer
values $\exp_{G_Q}f(nu)=[n]R$ are algebraic, but they do
not give an algebraic relation holding on the entire graph.
There is no contradiction with its full dimension.
One would need an additional arithmetic specialization
argument, or a way to obtain a further functional relation
from that single metric value. Neither is supplied by the
checked functional theorem. The endomorphism calculation
above also prevents obtaining extra CM multiplication
directions inside this fixed $G_Q$ by lifting the
noninteger endomorphisms of $E$.


\subsection{A CM-stable enlargement and its Gaussian eigencoordinate}

The fixed-extension calculation in \ref{ari:C.7.9} does not prevent
enlarging the torus. We test that alternative explicitly.
Retain its curve, its point $Q=2P$, and its real logarithm
$u$. Write $\beta=\sqrt{-19}=i\sqrt{19}$, $d=19$,
$D=\zeta(u)-s_2u$, and $k=\pi/A$. All endomorphisms and
extensions in what follows are defined over $\overline\Q$.

\begin{proposition}[the enlarged CM action does not isolate the Gaussian]\label{ari:C.7.10}
Let
\[
 \mathcal G=G_Q\mathbin{\times_E}G_{\beta Q},\qquad
 1\longrightarrow\mathbb G_m^2\longrightarrow\mathcal G
       \longrightarrow E\longrightarrow0.
\]
There is an endomorphism $B$ of $\mathcal G$ inducing
$\beta$ on $E$ and
$(t_1,t_2)\mapsto(t_2^{-1},t_1^d)$ on the torus.
It satisfies $B^2=[-d]$. Moreover
\[
 \operatorname{Hom}(\mathcal G,\mathbb G_m)=0,
 \qquad
 \operatorname{End}(\mathcal G)\otimes\Q=\Q(\beta).
\]
Thus $\mathcal G$ has no algebraic $\mathbb G_m$ quotient,
including after isogeny.

There is an algebraic compact point in the second extension
above $Q$ whose logarithm is
\[
 (u,v_2),\qquad v_2=-\beta uD+au,
 \qquad
 a=\frac1{2\beta}
       \frac{d}{du}\log F_\beta(Q)\in\overline\Q.
\]
Here the derivative means differentiation of the rational
function $F_\beta(\exp_E z)$ with respect to the algebraic
invariant differential, evaluated at $z=u$.
If $N(P)$ is algebraic, \ref{ari:C.7.9} combines this point with an
algebraic compact point in the first extension to give
\[
 U=(u,v_1,v_2),\qquad v_1=uD-ku^2,
 \qquad \exp_{\mathcal G}U\in\mathcal G(\overline\Q).
\]
In these algebraic logarithm coordinates the differential of
$B$ and its $-\beta$ eigenprojection are
\[
 dB=\begin{pmatrix}
 \beta&0&0\\ a&0&-1\\ \beta a&d&0
 \end{pmatrix},\qquad
 \Pi_-:=\frac12\left(I-\frac{dB}{\beta}\right),
\]
\[
 \Pi_-U=-\frac{\pi u^2}{2A}(0,1,\beta).
\]
The first torus coordinate of $\exp(\Pi_-U)$ is exactly
the Gaussian $\exp(-\pi u^2/(2A))$. The projector is not
the differential of an algebraic quasi-endomorphism. No
sequence of rational polynomials in $B$, even allowing
inverses of isogenies, can converge to it as complex linear
operators. Algebraicity of $\exp(U)$ does not supply
algebraicity of $\exp(\Pi_-U)$.
\end{proposition}
\Uses{the hypotheses, notation and definitions stated here; \cref{ari:C.7.9}; \cref{ari:C.7.6}; \cref{ari:eq:ce-zeta}}

\begin{proof}
Use the principal polarization to identify the two extension
classes with $Q$ and $\beta Q$. They are linearly independent
over $\Q$ modulo torsion: any nonzero $m+n\beta$ is an
isogeny, so $(m+n\beta)Q$ cannot be torsion. A character
of $\mathbb G_m^2$ given by $(m,n)\in\Z^2$ extends to
$\mathcal G$ exactly when the pushout extension has zero
class, namely when $mQ+n\beta Q=O$. Hence the restriction
of every homomorphism $\mathcal G\to\mathbb G_m$ to
the torus is zero. Such a homomorphism then factors through
$E$ and is zero. This proves the asserted absence of
characters. An isogenous semiabelian variety cannot acquire
a nonzero character, because composing it with the isogeny
would give a nonzero character of $\mathcal G$.

The Rosati adjoint of $\beta$ is $-\beta$. Pullback by
$\beta$ sends the ordered extension classes to
$(-\beta Q,dQ)$, exactly as does pushout by the displayed
torus matrix. Equality of these extension classes provides
the endomorphism $B$. It is unique because
$\operatorname{Hom}(E,\mathbb G_m^2)=0$. Its square
induces $[-d]$ on the torus and on $E$, so the same
uniqueness gives $B^2=[-d]$. The map
\[
 \operatorname{End}(\mathcal G)\otimes\Q
     \longrightarrow\operatorname{End}(E)\otimes\Q
          =\Q(\beta)
\]
is injective: an endomorphism inducing zero on $E$ has image
in the torus, and the absence of characters makes it zero.
Its image contains $\beta$, hence equals $\Q(\beta)$.

Next construct the compact point without assuming anything
about $N(P)$. The identities in \ref{ari:C.7.6}, applied at $\beta$
and $1+\beta$, give the algebraic nonzero number
\begin{equation}\label{ari:eq:ce-ribet-r}
 r_2=\frac{\theta((1+\beta)u)}{\theta(u)\theta(\beta u)},
 \qquad r_2^2=\frac{F_{1+\beta}(Q)}{F_\beta(Q)}.
\end{equation}
Indeed $(1+\beta)(1+\bar\beta)=1+d$, so the powers
of $\theta(u)$ cancel. None of the arguments is a lattice
point since $Q$ is nontorsion. Written in sigma functions,
\[
 r_2=\frac{\sigma((1+\beta)u)}{\sigma(u)\sigma(\beta u)}
                         e^{-s_2\beta u^2}.
\]
The point $(Q,r_2)$ in the rational trivialization of
$G_{\beta Q}$ therefore has logarithm
\[
 (u,v_2),\qquad
 v_2=u\{\zeta(\beta u)-s_2\beta u\}.
\]
Differentiating \ref{ari:C.7.6} shows that
\begin{equation}\label{ari:eq:ce-zeta}
 \zeta(\beta u)-s_2\beta u=-\beta D+a,
 \qquad a\in\overline\Q,
\end{equation}
with the value of $a$ in the statement. The compactness
character of this extension vanishes, since
\[
 \Re\{v_2-\zeta(\beta u)u+\beta u\eta(u)\}
    =\Re\left(\frac{\beta\pi|u|^2}{A}\right)=0.
\]
This also gives a direct, fully normalized construction of
the antisymmetric compact lift; no unrestricted theta
criterion is being used.

For completeness, the matrix $dB$ can be recovered without
choosing a transcendental change of coordinates. Its torus
block is the stated integer matrix and its first entry is
$\beta$, so initially write its first column as
$(\beta,a_1,a_2)^t$, with $a_1,a_2$ algebraic.
The identity $B^2=[-d]$ gives $a_2=\beta a_1$.
For $\omega\in\Lambda$, set
\[
 \kappa_1(\omega)=\zeta(u)\omega-\eta(\omega)u,
 \quad
 \kappa_2(\omega)=\zeta(\beta u)\omega-\eta(\omega)\beta u.
\]
Equation \eqref{ari:eq:ce-zeta} and the real-linear formula for
$\eta$ imply
\[
 \kappa_1(\beta\omega)=a\omega-\kappa_2(\omega).
\]
Comparing the second coordinates of the corresponding period
vectors under $B$ gives
$(a_1-a)\omega_j\in2\pi i\Z$ for $j=1,2$.
Since $\omega_2/\omega_1$ is nonreal, both integers must
be zero, so $a_1=a$. This proves the matrix formula.

Under the hypothetical algebraicity of $N(P)$, \ref{ari:C.7.9} gives
the first compact logarithm $v_1=uD-ku^2$. Pairing the
two algebraic lifts over the same base point gives $U$.
Substitution into the displayed matrix yields
\[
 v_2+\beta v_1-au=-\beta ku^2,
 \qquad
 \Pi_-U=-\frac{ku^2}{2}(0,1,\beta),
\]
which proves the exact Gaussian identity.

The two eigenspaces of $dB$ have eigenvalues $\beta$ and
$-\beta$. Every algebraic quasi-endomorphism has differential
$a_0I+b_0dB$ for rational $a_0,b_0$, with eigenvalues
$a_0+b_0\beta$ and $a_0-b_0\beta$. The projector
$\Pi_-$ would require these values to be $0$ and $1$,
which is impossible. There is also a uniform obstruction to
approximation. For arbitrary real $a_0,b_0$, put
$x=a_0+b_0\beta$, so that its other eigenvalue is
$\bar x$. Then
\begin{equation}\label{ari:eq:ce-gap}
 \max\{|x|,|\bar x-1|\}\ge\frac12,
\end{equation}
because $1\le|\bar x|+|\bar x-1|=|x|+|\bar x-1|$.
These are exactly the two eigenvalue errors in approximating
$\Pi_-$. The spectral radius is bounded by every induced
operator norm, so no such sequence can converge to the
projector. Compositions and inverses of isogenies still belong
to the same endomorphism algebra and satisfy this bound.

Finally the image of $\Pi_-$ is the toric Lie line
$\C(0,1,\beta)$. It is not the Lie algebra of an algebraic
subtorus: an algebraic one-dimensional subtorus of
$\mathbb G_m^2$ has rational slope. Since $u\ne0$,
the projected vector is nonzero. The analytic subgroup
theorem in fact shows that its exponential cannot be an
algebraic point of both torus coordinates. If the projection
preserved algebraic exponentials here, that would give a
contradiction; it is precisely the extra implication that
is unavailable. The construction isolates the Gaussian in
a Lie eigencoordinate, without producing an algebraic
group map that evaluates it.
\end{proof}

The spectral comparison with the ordinary prime $5$ is
explicit. The CM number $\alpha=(1+\beta)/2$ satisfies
$\alpha^2-\alpha+5=0$. Its two archimedean conjugates
both have absolute value $\sqrt5$. At the two places over
$5$, their normalized valuations are $0$ and $1$, because
the polynomial reduces to $X(X-1)$ and has product $5$.
The separation available in that nonarchimedean linear
algebra is absent in \eqref{ari:eq:ce-gap}. This comparison
does not turn a limiting operation on complex Lie spaces
into an algebraic-group endomorphism or transfer a
$p$-adic height conclusion to the real Gaussian.

\paragraph{An exact compact witness on the conductor-$361$ curve.}
In the preceding normalization the constants and the second
fibre coordinate are explicitly
\begin{equation}\label{ari:eq:ce-explicit-witness}
 s_2=-2,\qquad a=-\frac{21\beta}{703},\qquad
 r_2=-\frac{165}{74}+\frac{1839\beta}{1406}.
\end{equation}
We give polynomial certificates for these values, so they
do not depend on recognition from numerical approximations.
Put $f=4x^3-152x+361$ and
\[
 \begin{aligned}
 h={}&x^9-38x^8+437x^7-1444x^6-7942x^5\\
     &+82308x^4-274360x^3+390963x^2-130321x-130321,\\
 N={}&(-19x+76)h^2+f(h''h-(h')^2)+(6x^2-76)h'h,\\
 K={}&N'h-2Nh'.
 \end{aligned}
\]
Direct expansion gives
\[
 -K^2f=4N^3-152\cdot19^2Nh^4+361\cdot19^3h^6,
 \qquad \gcd(N,h)=1.
\]
Hence, in coordinates $Y=2y+1$, the map
\[
 (x,Y)\longmapsto
 \left(\frac{N}{19h^2},\frac{KY}{19\beta h^3}\right)
\]
is a degree-$19$ endomorphism fixing $O$, with pullback of
$dx/Y$ equal to $\beta\,dx/Y$. It is therefore the CM
endomorphism $\beta$. Divisors and the leading coefficient
at $O$ give $F_\beta=\beta^2h^2$. Its Laurent coefficient
of relative degree two is $-76$. Comparison with \ref{ari:C.7.6}
gives $38s_2=-76$, proving $s_2=-2$.

For the other kernel put $h_+=(H_0+\beta H_1)/10$, where
\[
 \begin{aligned}
 H_0={}&10x^8-380x^7+4864x^6-28880x^5+72561x^4\\
       &+54872x^3-795644x^2+1824494x-1433531,\\
 H_1={}&20x^7-304x^6+760x^5+10469x^4-83752x^3\\
       &+246924x^2-315514x+130321.
 \end{aligned}
\]
Writing $\bar h_+$ for replacement of $\beta$ by $-\beta$,
the finite polynomial identities and coprimality checks are
\[
 \begin{gathered}
 19xh^2-N=5f h_+\bar h_+,\qquad
 h_+\mid K+19\beta h^3,\\
 \gcd(h_+,fh)=\gcd(h_+,h_+')=1.
 \end{gathered}
\]
At a root of $h_+$ the displayed CM map sends $(x,Y)$ to
$(x,-Y)$, so these are precisely the eight pairs of
non-$2$-torsion points in $\ker(1+\beta)$. The other three
nonzero kernel points are $E[2]\setminus\{O\}$, since
$1+\beta=2\alpha$ in the CM order. Its norm is $20$,
which verifies the count. Thus
\[
 F_{1+\beta}=\frac{(1+\beta)^2}{4}f h_+^2.
\]
Taking the square ratio in \eqref{ari:eq:ce-ribet-r} and fixing
the sign by its expansion at $u=0$ gives the exact rational
function
\[
 r_2=-\frac{1+\beta}{2\beta}\frac{\wp'(u)h_+(\wp(u))}
                                             {h(\wp(u))}.
\]
At $Q$, the needed evaluations are
\[
 h(4)=-37,\qquad h'(4)=7,\qquad
 h_+(4)=\frac{829-279\beta}{10},\qquad\wp'(u)=-3.
\]
They give exactly \eqref{ari:eq:ce-explicit-witness}. In
particular $(Q,r_2)$ is an explicit nontorsion algebraic
point in the maximal compact subgroup of $G_{\beta Q}$.
Its torus-coordinate norm is
$r_2\bar r_2=974799/26011=4-\wp(\beta u)$.

As an independent numerical check, evaluation at $110$
decimal digits gives
\[
 \begin{aligned}
 c&=1.14793408195385447403276278660663356664777041472753\ldots,\\
 e^{-\pi u^2/(2A)}
  &=0.68133647525194864004469591167703191986051023921164\ldots.
 \end{aligned}
\]
The sigma quotient and the zeta expression for $r_2$ and $a$
agree with their exact values to more than $110$ digits.
The exact polynomial identities, evaluations and numerical
outputs are retained in \path{code/C7_cm_explicit.gp} and
its output. The polynomial construction proves the values;
the numerical comparison only checks the analytic branch.


\subsection{Complexifying the Gaussian: values, jets, and the rank criterion}

The antiholomorphic Gaussian can be made holomorphic by introducing a
second variable. This gives more algebraic values under the hypothesis
in C.7(c), but the derivative and growth hypotheses of the resulting
transcendence arguments must still be checked. We do this on the
conductor-$361$ model, where all constants and conjugations are fixed.

\begin{proposition}[the complexified route and its exact obstructions]\label{ari:C.7.11}
Use the curve, point and real logarithm of \ref{ari:C.7.9}:
\[
 E:y^2+y=x^3-38x+90,\qquad Q=(4,-2)=\exp_Eu,
 \qquad 0<u<\omega_1/2.
\]
Put $\beta=i\sqrt{19}$, $k=\pi/A$, $s_2=-2$, and
\[
 \theta(z)=\sigma(z)e^{-s_2z^2/2},\qquad
 D(z)=\zeta(z)-s_2z,\qquad
 \mathcal F(z,w)=\sqrt{19}\,\theta(z)\theta(w)e^{-kzw},
 \qquad H=\mathcal F^2.
\]
Then $\mathcal F(z,\bar z)=N(z)^2$, and $\mathcal F$ is an
entire function of two complex variables of order at most two.
It is periodic under $(z,w)\mapsto(z+\omega,w+\bar\omega)$,
$\omega\in\Lambda$. The subgroup
\[
 \Gamma=\{(\alpha u+\omega,\bar\alpha u+\bar\omega):
                   \alpha\in\Z[\beta],\ \omega\in\Lambda\}
 \subset\C^2
\]
has rank four over $\Z$ and lies in the real plane
$\{(z,\bar z):z\in\C\}$. If $N(Q)$ is algebraic, then the
values of $H$ on $\Gamma$ lie in one number field.

For every nonzero $\alpha\in\Z[\beta]$ there is an algebraic
number $a_\alpha$ such that
\[
 D(\alpha u)=\bar\alpha D(u)+a_\alpha.
\]
At $p_\alpha=(\alpha u,\bar\alpha u)$ the algebraic constant
direction $V_\alpha=\alpha\partial_z-\bar\alpha\partial_w$
satisfies
\begin{align*}
 V_\alpha\log\mathcal F(p_\alpha)
    &=\alpha a_\alpha-\bar\alpha\bar a_\alpha
                         \in\overline{\Q},\\
 V_\alpha^2\log\mathcal F(p_\alpha)
    &=-\alpha^2\{\wp(\alpha u)+s_2\}
      -\bar\alpha^2\{\wp(\bar\alpha u)+s_2\}
      +2\alpha\bar\alpha k.
\end{align*}
The number $k$ is transcendental. Consequently the second
directional derivative of $H$ at $p_\alpha$ is transcendental
whenever $H(p_\alpha)$ is algebraic. In particular,
\[
 (\partial_z-\partial_w)H(u,u)=0,\qquad
 \frac{(\partial_z-\partial_w)^2H(u,u)}{H(u,u)}=4k-8.
\]
Thus the cancelling first derivative does not give algebraic jets
of all orders. Moreover the rank-four set does not satisfy the
explicit derivative-free finite-order criterion audited below.
These are obstructions to the stated applications of those
criteria; they do not decide the algebraicity of $N(Q)$.
\end{proposition}
\Uses{the hypotheses, notation and definitions stated here; \cref{ari:C.7.9}; \cref{ari:C.7.1}; \cref{ari:C.7.6}; \cref{ari:eq:cc-fixedfield}; \citep[Theorem 6.1]{huber2022}; \cref{ari:C.7.7}}


\begin{proof}
The real model and its algebraic differential give real sigma
coefficients and a conjugation-stable lattice. In particular
$\overline{\theta(z)}=\theta(\bar z)$. The discriminant is
$-19^3$, so \ref{ari:C.7.1} gives the asserted diagonal identity.
The sigma function and the exponential of the quadratic
polynomial $-kzw$ have order at most two. From \ref{ari:C.7.6},
\[
 \theta(z+\omega)=\epsilon(\omega)
         e^{k\bar\omega(z+\omega/2)}\theta(z),
 \qquad \epsilon(\bar\omega)=\epsilon(\omega)\in\{1,-1\}.
\]
Multiplying this formula and its conjugate version, the two
signs multiply to one and the exponent cancels the change of
$-kzw$. This proves periodicity on all of $\C^2$, not only
on the real diagonal.

If $\alpha u+\omega=0$, then $[\alpha]Q=O$.
Nontorsion of $Q$ and the fact that a nonzero elliptic
endomorphism is an isogeny imply $\alpha=0$, and then
$\omega=0$. Thus the rank-two module $\Z[\beta]u$ and the
rank-two period lattice are independent over $\Z$.
This proves the rank assertion, with no extra algebraic
independence assumption on $u$ and the periods.

For $d=\alpha\bar\alpha$, the squared CM law \ref{ari:C.7.6} gives
the exact fixed-field identity
\begin{equation}\label{ari:eq:cc-fixedfield}
 H(\alpha u,\bar\alpha u)
   =19^{1-d}F_\alpha(Q)\overline{F_\alpha(Q)}H(u,u)^d.
\end{equation}
All the endomorphisms and their normalized rational functions
$F_\alpha$ are defined over a single number field: choose
a field containing the coefficients of the finitely many
generators of the CM order, and close it under conjugation.
The uniqueness of the Laurent normalization in \ref{ari:C.7.6} gives
the same field of definition for every $F_\alpha$.
Adjoining the hypothesized algebraic number $H(u,u)$ gives
one field for all the values in \eqref{ari:eq:cc-fixedfield}.
The value for $\alpha=0$ is zero, and period shifts do not
change it. Squaring was included because algebraicity of the
individual positive square roots alone would not establish
their membership in one fixed number field.

Logarithmically differentiating the squared CM law at $u$
gives
\[
 a_\alpha=\frac{1}{2\alpha}
     \frac{(d/dz)F_\alpha(\exp_Ez)|_{z=u}}{F_\alpha(Q)}
                  \in\overline{\Q}.
\]
The denominator is nonzero since $Q$ is nontorsion. The
derivation $d/dz$ is the algebraic invariant derivation on
$E$, so the displayed derivative has algebraic value.
Since $D(u)$ is real, conjugation gives the corresponding
formula at $\bar\alpha u$. On a neighborhood where
$\mathcal F\ne0$, its logarithmic gradient and Hessian are
\[
 (\partial_z\log\mathcal F,\partial_w\log\mathcal F)
       =(D(z)-kw,D(w)-kz),\qquad
 \operatorname{Hess}\log\mathcal F=
 \begin{pmatrix}-\wp(z)-s_2&-k\\-k&-\wp(w)-s_2\end{pmatrix}.
\]
Evaluation on $(\alpha,-\bar\alpha)$ proves both identities
in the statement. In particular the Gaussian cancels in the
first derivative and reappears in the second derivative.

To prove $k$ transcendental, recall the real period formula
$\eta_1=(s_2+k)\omega_1$. If $k$ were algebraic, the
nonzero period vector $(\omega_1,\eta_1)$ of the algebraic
universal vector extension of $E$ would lie on the algebraic
line of slope $s_2+k$. Dividing this period by a sufficiently
large integer supplies a nonzero algebraic torsion exponential.
The analytic subgroup theorem
\citep[Theorem 6.1]{huber2022} would give a one-dimensional
algebraic subgroup with that Lie line. Its projection to $E$
is nonzero, so it is an elliptic curve isogenous to $E$ and
splits the vector extension after that isogeny. This is
impossible: its nonzero class in $H^1(E,\mathcal O_E)$ stays
nonzero under an isogeny, whose pullback has a left inverse
after dividing the trace by the degree. These are the same
algebraic logarithm coordinates and extension as in \ref{ari:C.7.7};
a fixed algebraic linear change of coordinates leaves the
argument unchanged.

Write the two logarithmic derivatives as
$V_\alpha\log\mathcal F=B_\alpha$ and
$V_\alpha^2\log\mathcal F=A_\alpha+2dk$, where
$A_\alpha,B_\alpha$ are algebraic. Since $H=\mathcal F^2$,
\[
 \frac{V_\alpha^2H(p_\alpha)}{H(p_\alpha)}
       =4B_\alpha^2+2A_\alpha+4dk.
\]
The right side is transcendental, and the asserted conclusion
follows when $H(p_\alpha)\ne0$ is algebraic.
At $\alpha=1$, $a_1=0$, $\wp(u)=4$ and $s_2=-2$,
which give the last displayed formula in the statement.
\end{proof}

\paragraph{The precise several-variable criteria that were checked.}
The differential criterion of
\citet[Theorems 3.6 and 3.7]{waldschmidt2009auxiliary}
requires a finite algebraically independent family whose
partial derivatives belong to the polynomial ring generated
by that family over one number field. Algebraic values of
the full family then give algebraic values of all iterated
partial derivatives. The second-jet calculation above proves
that a family containing $H$ cannot satisfy both requirements
at the hypothesized algebraic orbit points. The issue persists
although the first cancelling direction is algebraic; the
direction also varies with $\alpha$.

For a criterion without differential equations we opened
\citet[Theorem 1, p.~24]{waldschmidt1975dimension}, including
its proof and a visual check of the displayed inequalities.
For $d$ meromorphic functions on $\C^n$ with growth and
arithmetic-size exponents $\rho_1,\ldots,\rho_d$, a set with
cardinality $N^{\ell-o(1)}$ in a ball of radius $O(N)$ and
separation at least $N^{-\lambda-o(1)}$ yields algebraic
dependence under the strict condition
\begin{equation}\label{ari:eq:cc-rankcriterion}
 (d-1)\ell>\sum_{j=1}^d\rho_j
                     +2d(n-1)(\lambda+1).
\end{equation}
The theorem also requires all function values to belong to
one number field and bounds on denominators in the entire
representations of the meromorphic functions. We do not
silently discard these hypotheses.

The natural independent triple here is
$\wp(z),\wp(w),H(z,w)$, all of order at most two.
The first two are algebraically independent. If $H$ satisfied
a polynomial over $\C(\wp(z),\wp(w))$, fix a generic pair
$(z,w)$ and translate only $z$ by $m\omega$, with
$0\ne\omega\in\Lambda$. The coefficients of the polynomial
remain constant, whereas the absolute values of its proposed
roots $H(z+m\omega,w)$ grow as
$\exp(k|\omega|^2m^2+O(|m|))$. This contradicts the finite
set of roots of a nonzero polynomial. Thus the triple is
independent.

Even granting all the size and denominator bounds with the
favorable exponents $\rho_j=2$, substitution of
$n=2$, $d=3$ and $\ell=4$ into
\eqref{ari:eq:cc-rankcriterion} would require
\[
 8>12+6\lambda,
\]
which fails even for $\lambda=0$. The actual geometry is
less favorable: the $\asymp N^4$ points obtained by bounding
the four integer coefficients lie in a ball of radius $O(N)$
in a real plane. Packing separated discs in that plane gives
$N^4\ll N^{2+2\lambda+o(1)}$, hence $\lambda\ge1$.
Removing the rank-two set of poles does not change this count.
Larger growth exponents only worsen the inequality. This
checks the particular published criterion rather than asserting
that every possible multivariable transcendence method fails.
The missing unconditional step in C.7(c) remains the exclusion
of an algebraic value on the diagonal compact leaf.


\subsection{Division values and a variable-degree interpolation test}

\begin{proposition}[route test, with an explicit limitation]\label{ari:C.7.12}
Let $E:y^2+y=x^3-38x+90$, $Q=(4,-2)$, and suppose, only in
this statement, that $N(Q)$ is algebraic.  Write
\[
 \delta=19^{1/4},\qquad b=\eta_1/\omega_1,\qquad
 g(z)=\delta\sigma(z)e^{-bz^2/2},\qquad Q=\exp_E(u),
\]
where $u,\omega_1$ are real as in \ref{ari:C.7.9}.  Thus
$|g(z)|=N(\exp_E z)$ for real $z$.

Every point in the CM division hull of $Q$ has algebraic norm.
There are constants $D,H$, depending on the fixed curve, $Q$ and
the hypothetical algebraic number $N(Q)$, such that, for
$nR=mQ$, $|m|\leq n$, $R\ne O$,
\[
 h(N(R))\leq H,
 \qquad
 [\Q(R,N(R)):\Q]\leq
 \begin{cases}Dn^4,&R\in E(\R),\\2Dn^6,&R\in E(\C).
 \end{cases}
\]
Here $h$ is absolute logarithmic height; the complex line means
the chosen complex embedding of an algebraic point.

Under the same hypothesis, any finite collection of real translates
$g(z+a)$ by division-hull logarithms has function-field
transcendence degree at most three over $\overline\Q$, after
adjoining the elliptic coordinates.  Three translates already
attain this cap.  The analogous holomorphic extensions in one
variable of $N(\exp_E(z+a))^2$ for complex division-hull shifts
have the same cap.  A safe common degree envelope for a fixed
finite family on the real division grids is $O(n^5)$, with
uniformly bounded heights.

These proved degree and height envelopes do not verify the
variable-degree criterion of Gramain--Mignotte--Waldschmidt
\citep[Theorem 2.1]{gramain1986values}.  This is a failure of the
available sufficient certificate, not a proof that the actual
degrees violate every possible application of that theorem.
The arithmetic-germ hypotheses in the criteria of
\citet[Theorem 6.1]{herblot2012meromorphic} and
\citet[Theorem 4.16]{gasbarri2010analytic} also remain unavailable;
the relevant nontorsion germs have a transcendental tangent.
\end{proposition}
\Uses{the hypotheses, notation and definitions stated here; \cref{ari:C.7.9}; \citep[Theorem 2.1]{gramain1986values}; \citet[Theorem 6.1]{herblot2012meromorphic}; \citet[Theorem 4.16]{gasbarri2010analytic}; \cref{ari:C.4.2}; \cref{ari:C.7.6}; \citep[Lemma 1.0]{gramain1986values}; \cref{ari:C.7.1}}


\begin{proof}
\emph{Propagation and degrees.}
The full division identity in \ref{ari:C.4.2} is
\[
 N(R)^{n^2}=N(nR)\delta^{n^2-1}/|\psi_n(R)|.
\]
For nontorsion $R$, its last factor is a nonzero algebraic number.
Consequently algebraicity propagates through multiplication
and division.  The CM multiplication formula in \ref{ari:C.7.6} gives
the same conclusion for the CM division hull.  All nonzero
torsion values were already proved algebraic in C.4.
For example, if $T$ has order dividing $n$ and $P$ is algebraic
and nontorsion, division at $P$ and $P+T$ gives directly
\[
 \left(\frac{N(P+T)}{N(P)}\right)^{n^2}
       =\frac{|\psi_n(P)|}{|\psi_n(P+T)|}\in\overline\Q.
\]
No undetermined harmonic normalization is needed.

Put $K_0=\Q(\delta,N(Q))$ and $D=[K_0:\Q]$.
For every nonzero integer $m$, $N(mQ)\in K_0$, since
$\psi_m(Q)$ is rational and its absolute value is rational.
The fibre of $[n]$ over $mQ$ has $n^2$ points, so
$[K_0(R):K_0]\leq n^2$.
If $R$ is real, $|\psi_n(R)|=\pm\psi_n(R)$; adjoining
$N(R)$ therefore costs at most a further factor $n^2$.
For a complex $R$, adjoin both $R$ and $\bar R$, at cost
at most $n^4$, and use
\[
 N(R)^{2n^2}=
 \frac{N(mQ)^2\delta^{2(n^2-1)}}
      {\psi_n(R)\psi_n(\bar R)}.
\]
This proves the stated degree envelopes when $m\ne0$.
For $m=0$, take the limit in the division identity at a nonzero
$n$-torsion point.  The zero of $\psi_n$ is simple in the
uniformizing coordinate, and $N(z)=\delta|z|+o(|z|)$ at zero;
hence
\[
 N(R)^{n^2}=
       \frac{n\delta^{n^2}}{|\partial_z\psi_n(R)|}.
\]
The derivative is algebraic, since the invariant derivation on
the Weierstrass model is algebraic.  The same degree bounds follow.

\emph{Uniform heights.}
We give the coefficient estimate needed here, so the height
assertion does not rely on an unspecified division-polynomial
bound.  On the fixed integral model, $\psi_n$ is a polynomial
in $x$ for odd $n$, and $(2y+1)$ times a polynomial in $x$
for even $n$.  These polynomials have degree $O(n^2)$.
Choose a fixed circle $|x|=A>1$ avoiding the finite branch values.
On its compact inverse image in $E(\C)$, $N$ has a positive
minimum $m_A$, and $|2y+1|$ has a positive minimum.  Let $M$
be the maximum of $N$ on the whole compact elliptic curve.
The rearranged division identity gives
\[
 |\psi_n(R)|\leq M\delta^{n^2-1}m_A^{-n^2}
       \qquad(|x(R)|=A).
\]
Cauchy's coefficient estimate, and division by $2y+1$ for
even $n$, show that the logarithm of the coefficient
$\ell^1$ norm is $O_E(n^2)$.  The same assertion holds for
$\partial_z\psi_n$, because differentiation and the fixed
Weierstrass equation add only polynomial factors in $n$.

The elementary height estimate for a fixed degree-four
duplication map, followed by telescoping, gives
$h(x(R))=2\widehat h(R)+O_E(1)$ uniformly over
$E(\overline\Q)$.  Since
$\widehat h(R)=(m/n)^2\widehat h(Q)$, the coordinate heights
are uniformly bounded when $|m|\leq n$.
The polynomial evaluation height inequality
\citep[Lemma 1.0]{gramain1986values} now gives
$h(\psi_n(R))=O_{E,Q}(n^2)$, and the same estimate for its
invariant derivative at torsion points.  Also
\[
 h(N(mQ))\leq m^2h(N(Q))+O_{E,Q}(m^2).
\]
Divide the height of the division identity by $n^2$.
For a complex value one uses
$h(|\alpha|)\leq h(\alpha)$, which follows by squaring and
using $|\alpha|^2=\alpha\bar\alpha$.
This proves a uniform bound for $h(N(R))$.

\emph{The translate count is finite in function-field dimension.}
For a real logarithm $a$ of a nonzero division point define
\[
 h_a(z)=\frac{\sigma(z+a)}{\sigma(z)\sigma(a)}e^{-baz}.
\]
Then
\[
 g(z+a)=\delta^{-1}g(z)g(a)h_a(z),\qquad
 h_a(z+\omega_1)=h_a(z),\quad
 h_a(z+\omega_2)=e^{-2\pi ia/\omega_1}h_a(z).
\]
The last two identities are exact sigma automorphy and Legendre's
relation.  A real point in the CM division hull is in the ordinary
division hull: if $nR=\alpha Q$ is real, then
$(\alpha-\bar\alpha)Q=O$, so $\alpha=\bar\alpha\in\Z$.
Indeed a nonzero endomorphism has finite kernel, whereas $Q$
is nontorsion.  The real locus of this negative-discriminant
curve is connected, so every real point has a real logarithm.
Consequently its real division-hull logarithms are precisely
of the form $(mu+k\omega_1)/n$, and lie in
$\Q u+\Q\omega_1$.
For finitely many shifts choose a common denominator
$q$ and write $a_j=(m_ju+k_j\omega_1)/q$; put $v=u/q$.
Then $h_{a_j}/h_v^{m_j}$ is elliptic for
$\Lambda_q=\Z\omega_1+\Z q\omega_2$.
The curve $\C/\Lambda_q$ is isogenous to $E$ and has an
algebraic model with this differential normalization.
Its displayed elliptic functions have algebraic divisors.
Each therefore equals a rational function over $\overline\Q$
times a complex scalar.  Evaluate at a real division point
avoiding their finitely many zeros and poles.  The formula
for $h_a$ in terms of $g$ and the assumed algebraicity of all
real division values make that scalar algebraic too.
Consequently all these translates are algebraic over
\[
 \overline\Q(E)\bigl(g,h_v\bigr).
\]
This proves the upper cap three over $\overline\Q$, not merely
over $\C$.

The cap is attained.  The multiplier
$e^{-2\pi iv/\omega_1}$ is not a root of unity, since $Q$
is nontorsion.  Applying the shifts $z\mapsto z+m\omega_2$
to a hypothetical polynomial relation shows that $h_v$ is
transcendental over the elliptic function field: at a fixed
generic $z$, a nonzero polynomial cannot vanish at infinitely
many distinct powers of this multiplier.  The same shifts
give $|g(z+m\omega_2)|=\exp(c m^2+O(m))$ with $c>0$.
After clearing denominators, any nonzero highest coefficient in
a polynomial in $g$, with coefficients in the elliptic field and
$h_v$, stays bounded
away from zero on a subsequence of the irrational rotation.
The highest power of $g$ then dominates, a contradiction.
Thus $g,h_v$ and the elliptic coordinate have transcendence
degree three.  Finally sigma addition gives
\[
 \wp(z)=\wp(v)-
 \frac{\delta^2g(z+v)g(z-v)}{g(z)^2g(v)^2},
\]
so the three translates $g(z),g(z+v),g(z-v)$ already attain it.

For a complex division shift $a$, the holomorphic extension
from real $z$ of the squared norm is
\[
 F_a(z)=\frac{N(a)^2}{\delta^2}
                  g(z)^2h_a(z)h_{\bar a}(z).
\]
The multiplier of $h_ah_{\bar a}$ depends only on
$a+\bar a\in\Q u+\Q\omega_1$ for this real CM curve.
The preceding finite-index-lattice argument applies again;
evaluating $F_a$ at a real division point makes the remaining
scalar algebraic.  Thus complex shifts of the norm do not
provide arbitrarily many independent functions either.

For completeness, the useful joint degree envelope for the
three generators is $O(n^5)$.  The map
$z\mapsto(\exp_Ez,h_v(z))$ describes the complex analytic
subgroup $\exp_{G_V}(z,(\zeta(v)-bv)z)$ of the algebraic
extension $G_V$, $V=\exp_E(v)$, in the meromorphic algebraic
section coordinate used in \ref{ari:C.7.1}.  Its restriction to real
$z$ is a compact circle; no compactness or closedness of the
whole complex analytic subgroup is asserted.
Its value at $z=u$ is algebraic, because
$h_v(u)=\delta g(u+v)/(g(u)g(v))$.
Its real period is $\omega_1$.  Hence for
$z=(mu+k\omega_1)/n$ its value is an $n$-division point of
a fixed-field algebraic multiple of that point.
The multiplication map on $G_V$ has degree $n^3$: its kernel
is an extension of $E[n]$ by $\mathbb G_m[n]$.
Adjoining $g(z)$ costs at most another $n^2$.
The fixed isogeny and the finitely many algebraic constants
affect only the implied constant.  Uniform height bounds for
the real translates follow from the same norm-bound proof with
denominator $nq$ and a fixed larger bound on the ratio of
numerator to denominator; quotient expressions give the same
bound for $h_v$.  For fixed complex shifts, use instead the
proved expression for $F_a$ as a rational elliptic function
on a fixed isogeny, times $g^2h_v^m$, with an algebraic scalar.
The heights of the isogenous elliptic coordinates, $g$ and
$h_v$ are bounded, so the rational-function evaluation height
inequality bounds $h(F_a)$ uniformly.  Removable zeros or poles
of this particular expression lie in a fixed finite set and
can be handled by their finitely many actual values.  Taking
the positive square root to recover the complex-shifted norm
divides its height by two and costs at most a fixed factor two
in degree for each of the fixed finite number of shifts.

\emph{The criterion has been tested, not assumed.}
For $k$ analytic functions, GMW Theorem 2.1 uses a point count
$\mu$, joint degree envelope $d$, height envelopes $\psi_j$
and growth envelopes $\theta_j$; its decisive sufficient
inequality is
\[
 d(N)\prod_{j=1}^k
  \max\{d(N)\psi_j(N),\theta_j(R(N))\}
 \leq c_0\mu(N)^{k-1}\varphi(N)^k,
\]
where
\[
 \varphi(N)=\log
 \frac{R(N)^2+r(N)r(N+1)}
      {R(N)(r(N)+r(N+1))}.
\]
We checked this formula and the other hypotheses in the
original printed pages 101--102.  On a fixed analytic disc
with fixed smaller sampling disc, $\varphi$ and the growth
envelopes are constant.  The height envelopes just proved
are constant as well.  Thus the required comparison reduces
to $d^{k+1}\ll\mu^{k-1}$.

At denominator $n$ there are $O(n^2)$ real points with
$nR=mQ$, $0\leq m\leq n$; through denominator $N$ there are
$O(N^3)$.  For $(\wp,g)$ on a disc without a pole, our proved
degree envelope $d=O(N^4)$ leads to the attempted comparison
$N^{12}\ll N^3$.  For three independent translates, the
safe joint envelope $d=O(N^5)$ leads to
$N^{20}\ll N^6$.  Neither comparison is available from
these estimates.  Adding further translates does not repair
this certificate: beyond three, the dependence conclusion
would already be true and would not contradict the norm
hypothesis.  We do not infer an actual degree lower bound
from an upper bound, and do not rule out a sharper arithmetic
estimate or a differently organized interpolation argument.

The alternative variable-degree theorems of Herblot and
Gasbarri require a formal germ over the residue number field,
with additional arithmetic estimates.  At a real nontorsion
division point with coordinate $z$ and algebraic $g(z)$,
\[
 \frac{dg}{dx}=\frac{g(z)}{\wp'(z)}
                         \bigl(\zeta(z)-bz\bigr).
\]
The slope in parentheses is transcendental by the period-line
argument of \ref{ari:C.7.9} applied to that nontorsion algebraic point.
The other two factors are nonzero algebraic numbers.
Thus the germ is not even defined over $\overline\Q$;
the required arithmetic-germ hypothesis has not been gained
by passing to division values.
\end{proof}

We also checked the Mahler-functional-equation route against
\citet[Theorem 1.1 and system (1.2)]{adamczewski2017mahler}.
That theorem requires algebraic Taylor coefficients, a linear
power-map functional system and an algebraic regular argument.
The compact normalization here has the transcendental coefficient
$b$, its elliptic logarithm is transcendental, and its duplication
in the algebraic coordinate is a Latt\`es map.  We have not obtained
a representation satisfying the stated hypotheses; we do not
assert that every alternative representation is impossible.
The entropy identity of
\citet[Theorem 5.1]{einsiedler2001entropy} uses a sequence of
transformations built from the denominators of $x(2^nQ)$,
with scale $4^n$.  It does not identify the height with the
entropy of a single finite-dimensional algebraic automorphism,
and supplies no algebraic value for its exponential.

The period-refined algebraic-group rank theorem has also been
checked in the original sources
\citep[Theorem 4.1 and Corollary 4.3]{waldschmidt1988periodrank}
\citep[main theorem, p.~363]{roy1993periodrank}.
Its coefficient is
$\alpha(G)=2\dim G-2d_0(G)-d_1(G)$, where $d_1$
counts multiplicative \emph{quotients}, not the dimension of
the torus subgroup.  For the nonsplit $G_Q$ of \ref{ari:C.7.9},
$d_0=d_1=0$: a nonzero character on its torus would kill
a nonzero integer multiple of the nontorsion extension class.
Thus $\alpha(G_Q)=4$.  The line $\ell_c$ has no nonzero
algebraic vector, has period rank $\kappa=1$, and the norm
hypothesis supplies one independent algebraic-point logarithm
in addition to that period.  The theorem allows total rank
$3=(4-1)\dim\ell_c$; the available rank $2$ therefore gives
no contradiction.  The period rank is exactly one because
a second period on this line would, by Legendre's relation,
force $u/\omega_1$ to be rational, contrary to the
nontorsion of $Q$.



\section{Limits and open points: the coefficient atlas}
\begin{remark}[what is and is not explicit]
\label{wb:rem:bcpr-literal-gauge}
Choose global $\bar\partial$-closed forms
\begin{equation}
 G\in\mathcal A^{0,1}(X),\qquad
 Q\in\mathcal A^{1,1}(X)
\end{equation}
representing $\mathfrak g$ and $\mathfrak q$.
\Cref{inv:A3.2} says that
$i_b^*[Q]_{\bar\partial}$ is the Dolbeault class of the explicit
translation-invariant
form displayed in \cref{wb:eq:bcpr-q-fibre-expanded}; it does not say that an arbitrary
representative $Q$ restricts literally to that form.  A literal global
coefficient formula would require additional $C^\infty$ gauge choices and
horizontal correction forms on the disc--pair-of-pants cover, compatible
with the two affine cyclic conjugacies and the toric cusp.  Those choices
are not part of the canonical data.  No such literal gauge is needed for
the class and product calculations below.
\end{remark}
\Uses{the hypotheses, notation and definitions stated here; \cref{inv:A3.2}; \cref{wb:eq:bcpr-q-fibre-expanded}}

\begin{remark}[literal global gauge boundary]
\label{wb:rem:bcpr-global-gauge-boundary}
The Bott--Chern algebra \cref{inv:C3.3} and the fibre-coordinate
restrictions displayed in
\cref{wb:eq:bcpr-gamma-form,wb:eq:bcpr-u-form,wb:eq:bcpr-q-fibre,wb:eq:bcpr-q-fibre-expanded,wb:eq:bcpr-q-gamma-product,wb:eq:bcpr-q-square}
are complete.  The canonical data do not select a preferred literal global
coefficient atlas for the smooth forms $G$ and $Q$.  Such an atlas would
require explicit smooth local trivializations and horizontal correction
functions through both finite affine fillings and the cusp.  After fixing
$\mathfrak g$ and one coherent-Leray lift $\mathfrak q$, changing those
smooth gauges changes $G$ and $Q$ by Dolbeault coboundaries but changes
none of the cohomology classes, nonvanishing arguments, or products proved
above.  This is a limitation only on a preferred global coefficient
presentation, not on the representatives as cohomology classes or on the
Bott--Chern algebra.
\end{remark}
\Uses{the hypotheses, notation and definitions stated here; \cref{inv:C3.3}; \cref{wb:eq:bcpr-gamma-form,wb:eq:bcpr-u-form,wb:eq:bcpr-q-fibre,wb:eq:bcpr-q-fibre-expanded,wb:eq:bcpr-q-gamma-product,wb:eq:bcpr-q-square}}


\begin{question}[Open: A literal global coefficient atlas for the generators]\label{inv:C3.8}
The invariant computations and all the metric types named above are
settled for $X$.  The real restriction map
\[
r_b:H_A^{1,1}(X;\R)\longrightarrow H_A^{1,1}(F_b;\R)
\]
has image $\R\eta$ by \ref{inv:A3.2} and \ref{inv:C3.7}.  Its source has real dimension
two and its target has real dimension four.  This concerns
cohomology classes and does not assert literal equality of arbitrary
smooth representatives with an invariant fibre form.

The formula \eqref{inv:eq:C-G-Green} supplies a global $G$ after a metric
choice.  For the standard primitive marking, $G_\gamma=G_0/12$
represents $g_\gamma$ by \ref{inv:C3.2}.  There is a precise operator
construction of $Q$ as well.  Fix a Hermitian metric on $X$ and a
regular fibre $F$, identify $\eta$ with its real invariant form on
$F$, and choose an orthonormal basis $e_1,e_2$ of the Dolbeault-harmonic
$(1,1)$ forms for $\Delta_{\bar\partial}$.  Define
\[
\ell(e)=-\frac1{12}\int_F i_F^*e\wedge\eta,\qquad
Q=\frac{\overline{\ell(e_1)}e_1+
            \overline{\ell(e_2)}e_2}
          {|\ell(e_1)|^2+|\ell(e_2)|^2}.
\]
The denominator is positive by \ref{inv:A3.2} and
$\int_F\eta^2=-12$.  The expression is independent of the
orthonormal basis, is $\bar\partial$-closed, and has fibre
Dolbeault class $\eta$.  Indeed the image of fibre restriction on
$H^{1,1}_{\bar\partial}(X)$ is $\C\eta$, with kernel
$\C d_1g$; the displayed functional measures exactly the coefficient
of $\eta$, and the formula makes that coefficient one.

These two formulas fix global smooth representatives by explicit
elliptic operators and a finite normalization.  We have not evaluated
their coefficient functions on a common atlas through the three
special fibres.  The remaining task is therefore precise: compute
the harmonic projector and the two displayed restriction integrals
for one fully specified global Hermitian metric, or construct and
evaluate a Cech representative and partition-of-unity correction
with the same restrictions.  No numerical computation of those
infinite-dimensional elliptic inverses was run.  This is the
unfulfilled literal-coefficient requirement, not an undecided Hodge
number, product, or existence of the Leray class.
\end{question}
\Uses{the hypotheses, notation and definitions stated here; \cref{inv:A3.2}; \cref{inv:C3.7}; \cref{inv:eq:C-G-Green}; \cref{inv:C3.2}}

\section{Unconditional arithmetic of the normalized norm}\label{ari:open:norm}

Statements \ref{ari:C.7.1}--\ref{ari:C.7.5} do not decide the unconditional
arithmetic nature of $N$ at a nontorsion point. An explicit
small instance is
\[
 E:y^2=x^3-140x+784,\qquad P=(0,28),\qquad j=-3375,
 \qquad N(P)=2e^{-\widehat h(P)}.
\]
Here $P$ is a primitive generator modulo the rational
$2$-torsion and
\[
 \widehat h(P)=
 0.46736567718147128000752020126284678363562034370331\ldots.
\]
The exact unresolved assertion is that $2e^{-\widehat h(P)}$
is transcendental. The corresponding example of conductor
$361$, $y^2+y=x^3-38x+90$, $P=(0,9)$, has
$N=e^{-\widehat h(P)}$ and no finite local contribution.
The relative-period proof isolates its obstruction: an
unconditional theorem excluding
\eqref{ari:eq:cp-relation} for the positive diagonal central
direction. A theorem about linear $1$-period relations
does not supply that exclusion. The attempted approaches and the exact CM theta comparison are printed in the preceding period and CM-normalization sections;
their source checks and computational records are retained in the audit files.
These arguments do not resolve the stated exclusion.

\clearpage
